\documentclass[10pt]{beamer}
\usetheme{Madrid}
\usepackage[T1]{fontenc}
\usepackage[utf8]{inputenc}
\usepackage[brazil]{babel}
\usepackage{lmodern}
\usepackage{hyperref}
\setbeamertemplate{navigation symbols}{}

\title{Seguidor de Linha}
\subtitle{Curso de Microcontroladores}
\author{}
\date{}

\begin{document}

\begin{frame}
  \titlepage
\end{frame}

\begin{frame}{Introdução}
  \begin{itemize}
    \item Um seguidor de linha é um robô que segue uma pista no chão.
    \item O robô lê a linha com sensores, processa a informação e ajusta os motores.
    \item O objetivo é permanecer sobre a linha e seguir seu trajeto.
    \item É um robô autônomo que se desloca sobre um caminho marcado.
    \item A pista é uma linha branca em fundo escuro.
    \item A ideia central é detectar a linha e corrigir a direção.
  \end{itemize}
\end{frame}

\begin{frame}{Ideia básica}
  \begin{itemize}
    \item Use um número ímpar de sensores.
    \item O sensor central deve sempre ler a linha branca.
    \item Se um sensor lateral ler a linha, o motor oposto deve ganhar mais velocidade.

    \item Pode ficar instável em curvas mais fechadas.
    \item Ruídos nos sensores podem gerar comandos opostos, porque a leitura é binária.
    \item O robô tende a fazer ajustes bruscos, não suaves.
  \end{itemize}
  \vspace{0.4cm}
  Por isso usamos um controle melhor, como PID.
\end{frame}

\begin{frame}{PID}
  PID é uma técnica de controle que melhora o ajuste do robô.

  \begin{itemize}
    \item \textbf{Proporcional (P)}: corrige com base no erro atual.
    \item \textbf{Integral (I)}: corrige com base no histórico do erro.
    \item \textbf{Derivativo (D)}: corrige com base na velocidade do erro.
  \end{itemize}

  \pause

  No seguidor de linha, o PID ajuda a fazer curvas suaves e manter o robô na pista.
\end{frame}

\begin{frame}{Divisão do seguidor}
  O seguidor de linha tem três partes principais e uma extra:
  \begin{itemize}
    \item Sensores
    \item Microcontrolador
    \item Motores
    \item PCB (parte extra)
  \end{itemize}
\end{frame}

\begin{frame}{Sensores}
  \begin{itemize}
    \item Os sensores são a entrada do robô.
    \item Eles detectam a linha no chão e enviam a informação ao micro.
    \item Sensores e ESP32 exigem conhecimento de microcontrolador e eletrônica.
  \end{itemize}
\end{frame}

\begin{frame}{Microcontrolador}
  \begin{itemize}
    \item O micro é onde tudo acontece.
    \item Ele lê os sensores, aplica a lógica e gera comandos para os motores.
    \item O ESP32 calcula o PID e decide a velocidade de cada roda.
  \end{itemize}
\end{frame}

\begin{frame}{Motores}
  \begin{itemize}
    \item Os motores são a saída do sistema.
    \item Eles transformam o sinal elétrico em movimento.
    \item A ponte H permite controlar direção e velocidade.
    \item Essa parte envolve eletrônica e mecânica.
  \end{itemize}
\end{frame}

\begin{frame}{Parte extra: PCB}
  \begin{itemize}
    \item A PCB organiza os componentes eletrônicos.
    \item Ela envolve mais design do que programação.
    \item Um bom projeto de PCB deixa a montagem mais fácil. E pode evitar que o seguidor seja muito frágil (tanto em termos eletrônicos quanto mecânicos)
  \end{itemize}
\end{frame}

\begin{frame}{Resumo}
  \begin{itemize}
    \item Sensores leem a pista.
    \item Micros processam e calculam os ajustes.
    \item Motores movem o robô.
    \item PCB organiza o hardware.
    \item PID melhora o controle e torna o seguidor mais suave.
  \end{itemize}
\end{frame}

\end{document}
